\chapter{枚举与模式匹配}
\label{ch:10}

枚举把一个值的几种状态写进类型。每个状态可以不带数据，也可以携带一个或多个载荷。match 根据状态选分支，并把载荷绑定为局部变量。

\section{状态与载荷}

先列出互斥状态，再决定各自的数据。例如请求可以是等待、完成或失败，完成状态携带正文，失败状态携带错误。读取载荷前必须先识别所在变体；匹配表达式需要覆盖可达状态。

\par
\medskip
\begin{lstlisting}[style=chali]
enum Packet
{
    Data(String text), Empty
};
void main()
{
    Packet packet = Packet::Data("original");
    String label = match (packet)
    {
        Packet::Data(text) =>
        {
            text = "copy";
            yield text;
        }
        Packet::Empty =>
        {
            "empty"
        }
    };
    String original = match (packet)
    {
        Packet::Data(text) =>
        {
            text
        }
        Packet::Empty =>
        {
            "empty"
        }
    };
    println("{} {}", label, original);
}
\end{lstlisting}
\par
\medskip

\par
\medskip
\begin{lstlisting}[style=out]
copy original
\end{lstlisting}
\par
\medskip

\section{多个载荷与通配}

多个载荷按声明顺序写入括号。普通载荷可以直接绑定变量；不需要的载荷写 \_。如果载荷是接口，类型细分使用 Type value 形式，外层枚举仍写出。第 13 章再说明接口值的具体类型匹配。

\par
\medskip
\begin{lstlisting}[style=chali]
enum Event
{
    Idle,
    Message(String sender, String text)
};
void main()
{
    let event = Event::Message("Ada", "ready");
    match (event)
    {
        Event::Idle => { println("idle"); }
        Event::Message(sender, text) => { println("{}: {}", sender, text); }
    }
    match (event)
    {
        Event::Idle => { println("idle"); }
        Event::Message(_, text) => { println("{}", text); }
    }
}
\end{lstlisting}
\par
\medskip

\par
\medskip
\begin{lstlisting}[style=out]
Ada: ready
ready
\end{lstlisting}
\par
\medskip

\section{复制与派生能力}

载荷沿用其类型的值、借用和资源规则。复制一个带 String 的枚举后，修改副本不会通过值语义改变原值；引用载荷则仍然表达借用关系。不要把枚举包装当成取消所有权检查的方法。

\par
\medskip
\begin{lstlisting}[style=chali]
#derive(Equal, Hash, Clone)
enum BookPacket
{
    Empty,
    Text(String text)
}

void main()
{
    BookPacket first = BookPacket::Text("value");
    BookPacket second = first.clone();
    assert(!(first == BookPacket::Empty));
    println("{} {}", first == second, first.hash(first) == second.hash(second));
}
\end{lstlisting}
\par
\medskip

\par
\medskip
\begin{lstlisting}[style=out]
1 1
\end{lstlisting}
\par
\medskip

\section{将 match 用于计算}

match 也可以产生值，各分支结果必须满足共同类型要求。避免在后面依靠断言猜测变体：把需要载荷的操作写进对应分支，使类型和控制流一起记录前提。Option 和 Result 是这种结构在缺省与错误处理中的应用。

\section{习题}

\begin{enumerate}

\item 给 Event 增加断开状态。

\item 只打印 Message 的正文。

\item 何时需要通配分支？

\item 引用载荷是否等于独立副本？

\end{enumerate}

\subsection*{参考解答与验收要点}

\begin{enumerate}

\item 同步更新所有需要穷尽的匹配。

\item 写 Message(\_, text) 并使用 text。

\item 未逐个覆盖剩余状态或接口实现时需要保留回退。

\item 不是；生命周期和访问限制仍然适用。

\end{enumerate}

下一章：\hyperref[ch:11]{第 11 章}。

\section{番外：大灰狼漏掉的分支}

改写自传统童话《三只小猪》。

从前有三只小猪，离家各自造屋。老大造了一间草屋，老二造了一间木屋，老三造了一间砖屋。大灰狼听见商机，连夜写了一份拆屋程序，号称支持一切建筑，实际只有两个分支。

大灰狼来到草屋前，说：“小猪，小猪，让我进来。”小猪说：“不让你进来！”狼于是深吸一口气，一吹，草屋倒了。测试通过。

它又来到木屋前，说：“小猪，小猪，让我进来。”小猪仍不让。狼按手册一推，木屋倒了。测试覆盖率百分之百——因为它只把自己会的两种房屋登记进了统计表。

到了砖屋，狼又吹又推，房子不动。它翻到兜底分支，那里写着：“其余按木屋处理。”第三只小猪隔窗问：“你这叫兼容砖屋，还是把砖屋改名成木屋？”

狼想从烟囱钻进去，跌进一口热锅。锅底贴着错误信息：“未知变体，拒绝执行。”

三只小猪没有再争论兜底分支短不短。它们把房屋分成清楚的变体，让每种结果都有去处。大灰狼后来改行做代码审查，见到 \texttt{\_\allowbreak{}} 就先问一句：“这个以后会不会装下一口锅？”

\subsection*{技术对照}

枚举变体表达有限分类。match 的覆盖与守卫需要一起考虑；兜底分支可能隐藏新增变体，重叠警告也不应被误写成所有场合都拒绝编译。
